\documentclass[conference]{IEEEtran}

\usepackage[utf8]{inputenc}
\usepackage[russian]{babel}
\usepackage{amsmath, amssymb}
\usepackage{graphicx}

\begin{document}

\title{SFLXGB: подход к восстановлению пропусков во временных рядах на основе сезонного лагового градиентного бустинга и ансамблевых методов импутации}

\author{\IEEEauthorblockN{Автор}
\IEEEauthorblockA{Организация}
}

\maketitle

\begin{abstract}
Качество прогнозирования временных рядов существенно зависит от обработки пропусков. В работе рассматриваются два метода восстановления данных: SFLXGB и AIM, направленные на сохранение динамической структуры временных рядов и повышение качества прогнозирования при наличии значительных пропусков.
\end{abstract}

\begin{IEEEkeywords}
временные ряды, импутация, градиентный бустинг, ансамблевые методы, прогнозирование
\end{IEEEkeywords}

\section{Введение}

Качество прогнозирования временных рядов напрямую зависит от корректности этапа подготовки данных. Одной из ключевых проблем является наличие пропусков, нарушающих временную структуру и искажающих динамические зависимости. Стандартные методы обработки часто недостаточно эффективны при сложных паттернах данных, особенно в условиях нестационарности и шумов.

В прикладных задачах временные ряды редко бывают полными. Пропуски возникают из-за сбоев датчиков, потерь телеметрии, ошибок передачи данных и регламентного обслуживания систем. Это приводит к неполным наблюдениям, снижает устойчивость моделей и ухудшает качество анализа.

Особенно критична данная проблема в энергетике и промышленном мониторинге, где прогнозирование нагрузок является основой управления инфраструктурой. В условиях цифровой трансформации возрастает потребность в интеллектуальных методах анализа, способных работать с неполными и быстро меняющимися данными.

Таким образом, задача восстановления пропусков является критически важной для повышения качества моделей прогнозирования.

\section{Математическая модель}

Временной ряд задается как:
\begin{equation}
y = (y_1, y_2, \dots, y_T)
\end{equation}

где $y_t$ — наблюдение в момент времени $t$.

Частично наблюдаемое состояние:
\begin{equation}
x_t^{partial} = (y_{t-1}, \dots, NaN, \dots, y_{t-p})
\end{equation}

Восстановление состояния:
\begin{equation}
\tilde{x}_t = R(x_t^{partial})
\end{equation}

Гипотеза эквивалентности:
\begin{equation}
F(x_t^{complete}) \approx F(R(x_t^{partial}))
\end{equation}

Локальное восстановление:
\begin{equation}
y_{t-i}' =
\begin{cases}
y_{t-i}, & \text{если наблюдается} \\
g(\cdot), & \text{если отсутствует}
\end{cases}
\end{equation}

Формирование полного вектора:
\begin{equation}
\tilde{x}_t = (y_{t-1}', y_{t-2}', \dots, y_{t-p}')
\end{equation}

\section{SFLXGB}

Оператор восстановления:
\begin{equation}
\tilde{x}_t = R(x_t^{partial})
\end{equation}

Гипотеза:
\begin{equation}
F(x_t^{complete}) \approx F(R(x_t^{partial}))
\end{equation}

Локальный оператор:
\begin{equation}
g(x_t^{lag}, t) =
\begin{cases}
f_\theta(x_t^{lag}), & \text{если нет пропусков} \\
g_{HDIRT}(y_{t-w:t-1}), & \text{иначе}
\end{cases}
\end{equation}

\section{AIM}

Класс пропуска:
\begin{equation}
c_t = \phi(l_t)
\end{equation}

Выбор метода:
\begin{equation}
m^*(c) = \arg\min_{m_k \in M} L_k^{(c)}
\end{equation}

Итоговое восстановление:
\begin{equation}
\hat{y}_t =
\begin{cases}
m^*(c_t)(D_{miss}), & t \in \Omega_{gap} \\
y_t, & t \notin \Omega_{gap}
\end{cases}
\end{equation}

\section{Программная апробация}

Использованы датасеты: Elista, Istanbul Traffic, Temperature.

Рассмотрены методы: SFLXGB, AIM, HDIRT, SKNN, KNN, LR, MEAN, MEDIAN, LAST, SMEAN, XGB.

Эксперименты проводились при уровнях пропусков 10\%, 30\%, 50\%, 70\%.

Формирование пропусков осуществлялось путем генерации случайных интервалов различной длины с последующим размещением без пересечений.

\end{document}